\cotefile{157-2}
\scanpage{323}
\mspage{62}
% archive sheet number: 322
% ms. 1–2
de type $\Delta_+$, ou aussi au \deleted{\unsure{n'importe quel}}\inserted{chemin dégénéré} de n'importe quel type $\tau$, disons
% ms. 3
\diagram{\begin{tikzcd}[ampersand replacement=\&,column sep=small]
x \& x\arrow[l,"\mathrm{id}"']\arrow[r,"\mathrm{id}"] \& x \arrow[r,dash,"\mathrm{id}"] \& x \& x\arrow[l,"\mathrm{id}"'] \& x\arrow[l,"\mathrm{id}"']{,}
\end{tikzcd}}
% ms. 4
un tel chemin n'étant autre que $p_\tau^*(x)$, où
% ms. 5
\[p_\tau:\tau\longrightarrow 1\]
% ms. 6–7
est le morphisme de \unsure{dégénérescence} (l'unique morphisme dans $\underline{\unsure{\mathrm{Ty}|\mathrm{Tz}}}$ vers $1$, qui est objet final),
% ms. 8–9
d'où
\diagram{\begin{tikzcd}[ampersand replacement=\&,row sep=0pt]
p_\tau^*:\underline{\mathrm{Ch}}_1(\deleted{\unsure{x}}x,x)=\{x\}\arrow[r] \& \underline{\mathrm{Ch}}_\tau(x,x)\\
\& \cap\\
\& \underline{\mathrm{Ch}}(x,x).
\end{tikzcd}}

% ms. 10–12
La catégorie \inserted{$\Pi(X)$} ainsi obtenue doit s'appeler le \uline{groupoïde fondamental} de $X$. En effet, \deleted{\unsure{a)}}

% ms. 12–15
a) toutes les flèches de $\Pi(X)$ sont inversibles, i.e. $\Pi(X)$ est bien un groupoïde, et

% ms. 16
\unsure{b)} le foncteur canonique
% ms. 17
\[X\longrightarrow\Pi(X)\]
% ms. 18–22
qui est l'identité sur les objets, et qui est défini sur les flèches par $\unsure{u\mapsto\mathrm{cl}\,u|u\mapsto\mathrm{cl}(u)}$ comme ci-dessus, est \uline{universel} \inserted{(dans $\mathrm{Cat}$)} pour les foncteurs \inserted{(i.e. les $\mathrm{Cat}$-morphismes)} de $X$ vers des \uline{groupoïdes}.
% ms. 23–26
(Plus généralement, il est universel pour les foncteurs $X\xrightarrow{F}Y$ tels que $F$ transforme \unsure{tte} flèche de $X$ en flèche inversible.)
\scanpage{324}
\mspage{63}
% archive sheet number: 323
% ms. 1–5
\deleted{\unsure{Comme}}Pour prouver que \unsure{tte} flèche dans $\deleted{\unsure{X}}\Pi(X)$ est inversible, comme \deleted{\unsure{tout}} \unsure{tt|$\varphi$} chemin \inserted{non vide} est \unsure{composé|comp. sé} de chemins ``élémentaires'' de type $\Delta_+$ ou $\Delta_-$, il suffit de prouver que ces chemins sont inversibles. En fait, ils se correspondent biunivoquement
% ms. 6
\[\underline{\mathrm{Ch}}_{\Delta_+}(a,b)\simeq\underline{\mathrm{Ch}}_{\Delta_-}(b,a)\]
% ms. 7
\deleted{\unsure{corresp.}} associant au chemin de $a$ vers $b$
% ms. 8
\[c(u):\quad a\xrightarrow{u}b\qquad\text{de type }\Delta_+\]
% ms. 9
le chemin
% ms. 10
\[c(u)^{\unsure{\circ|o}}:\quad b\xleftarrow{u}a\qquad\text{de type }\Delta_-\]
% ms. 11
de $b$ vers $a$. Je dis que
% ms. 12–13
\[\begin{aligned}
c(u)^\circ\,c(u)&\sim\deleted{\unsure{\mathrm{id}}}1_a\\
c(u)\,\unsure{(c(u))^\circ|c(u)^\circ}&\sim 1_b
\end{aligned}\]
% ms. 14–17
\diagram{\begin{tikzcd}[ampersand replacement=\&,column sep=small,row sep=small]
a\arrow[r,"u"] \& b \& a\arrow[l,"u"'] \& \\
a\arrow[u,"{\deleted{\unsure{p}}\,\mathrm{id}}" description]\arrow[r,"\mathrm{id}"] \& {\overset{\deleted{\mathrm{id}}}{a}} \arrow[u,"u"'] \& a\arrow[u,"\mathrm{id}"']\arrow[l,leftrightarrow,"\mathrm{id}"'] \& {\left.\vphantom{\begin{matrix}x\\x\\x\\x\end{matrix}}\right\}\ \begin{gathered}\text{on a un}\\ \text{hom de }\mathrm{Ch}_{\Delta_+\Delta_-}(a,a)\\ \text{de }1_a\text{ dans }c(u)^\circ c(u)\end{gathered}}\\
{}\arrow[u] \& a \& \&
\end{tikzcd}}
% ms. 18–21
\diagram{\begin{tikzcd}[ampersand replacement=\&,column sep=small,row sep=small]
b \& a\arrow[l,"u"']\arrow[r,"u"] \& b \& \\
a\arrow[u,"u"] \& a\arrow[l,"\mathrm{id}"']\arrow[u,"\mathrm{id}"]\arrow[r,"\mathrm{id}"] \& a\arrow[u,"u"'] \& {\left.\vphantom{\begin{matrix}x\\x\\x\\x\end{matrix}}\right\}\ \begin{gathered}\text{item dans }\mathrm{Ch}_{\Delta_-\Delta_+}(b,b)\\ 1_b\longrightarrow c(u)c(u)^\circ\end{gathered}}\\
{}\arrow[u] \& a \& \&
\end{tikzcd}}
\scanpage{325}
\mspage{64}
% archive sheet number: 324
% ms. 1–3
\deleted{\unsure{La démonstration de}} Notons que \inserted{plus \unsure{génért}}, pour tout chemin $c\unsure{;:|:}\,a\longrightarrow b$, \inserted{de type $\tau$,} on définit un \uline{chemin opposé}
% ms. 4
\[c^{\unsure{\circ|o}}:b\longrightarrow a\qquad\text{de type }\tau^\circ,\]
% ms. 5
où
% ms. 6
\[\tau\longmapsto\tau^\circ:\quad\mathrm{Ty}\longrightarrow\mathrm{Ty}\]
% ms. 7–10
est l'\unsure{anti-involution} \deleted{\unsure{de}} du monoïde libre $\mathrm{Ty}$ qui échange les deux générateurs $\Delta_-,\Delta_+$, et qui \deleted{\unsure{sur}} est donc donné sur l'élément général par
% ms. 11
\[(\Delta_{\varepsilon_1}^{\deleted{\unsure{\varepsilon_1}}}\,\Delta_{\varepsilon_2}^{\deleted{\unsure{\varepsilon_2}}}\,\cdots\,\Delta_{\varepsilon_n})^\circ
=\Delta_{-\varepsilon_n}\,\Delta_{-\varepsilon_{n-1}}\cdots\Delta_{-\varepsilon_{\unsure{n-|n|1}}}.\]
% ms. 12–15
Par récurrence \uline{sur la longueur} de $c$, en \unsure{utilisant} \deleted{\unsure{enfin}} le résultat précédent \deleted{que} que $c(u)$, $c(u)^\circ$ sont inverses l'un de l'autre dans $\Pi(X)$, on trouve
% ms. 12–15 left
% slanted
\margin{inutile, \unsure{résulte} directement de la formule $(*)$ page suivante, \unsure{en} décomposant $c$ en \unsure{composé|suite} de chemins ``\unsure{atomiques|élémentaires}'' de type $\Delta_+$ ou $\Delta_-$.}
% ms. 16–17
\[\left.\begin{aligned}c^\circ c&=\mathrm{id}_a\\cc^\circ&=\mathrm{id}_b\end{aligned}\right)\quad\text{dans }\Pi(X).\]

% ms. 18–19
\uline{Exemple} Si
% ms. 18–19
\[\begin{aligned}\tau&=\Delta_+\,\Delta_-\,\Delta_+^2\,\Delta_{\unsure{-|+}}\\
\tau'&=\Delta_+\,\Delta_-^2\,\Delta_+\,\Delta_-\end{aligned}\]
% ms. 20
Et si on a un chemin
% ms. 21
\diagram{\begin{tikzcd}[ampersand replacement=\&,column sep=small]
c:\quad a_0\arrow[r,"u_1"] \& a_1 \& a_2\arrow[l,"u_2"']\arrow[r,"u_3"] \& a_3\arrow[r,"u_4"] \& a_4 \& a_5\arrow[l,"u_5"']
\end{tikzcd}}
% ms. 22
de type $\tau$, on en déduit \deleted{\unsure{des}} le chemin
% ms. 23
\diagram{\begin{tikzcd}[ampersand replacement=\&,column sep=small]
\underline{c}^\circ:\quad a_5\arrow[r,"u_5"] \& a_4 \& a_3\arrow[l,"u_4"'] \& a_2\arrow[l,"u_3"']\arrow[r,"u_2"] \& a_1 \& a_0\arrow[l,"u_1"']
\end{tikzcd}}
% ms. 24
de type $\tau'$. L'opération $c\longmapsto c^\circ$ satisfait en
\scanpage{326}
\mspage{65}
% archive sheet number: 325
% ms. 1
effet
% ms. 2–3
\diagram{\begin{tikzcd}[ampersand replacement=\&,row sep=0pt,column sep=small]
a\arrow[r,"\gamma"] \& b\arrow[r,"\gamma'"] \& c\\
c\arrow[r,"{\gamma'^{\unsure{\circ|o}}}"] \& b\arrow[r,"\gamma^\circ"] \& a
\end{tikzcd}}
% ms. 4
la relation
% ms. 5
\[(\gamma'\unsure{\,.\,|\,}\gamma)^\circ=\gamma^\circ\gamma'^\circ,\tag*{$(*)$}\]
% ms. 6–9
ce qui est en accord avec l'intuition que l'opération $\gamma\longmapsto\gamma^\circ$ joue le rôle d'un passage à l'inverse (``à homotopie près''), et avec la formule
% ms. 9
\[(\gamma'\gamma)^{-1}=\gamma^{-1}\gamma'^{-1}\]
% ms. 10–11
pour deux flèches \deleted{\unsure{inv.}} $\gamma,\gamma'$ composables (dans l'ordre $\gamma',\gamma$).
% ms. between 11 and 12
\figure{\begin{tikzpicture}[x=1cm,y=1cm]\draw (0,.20) -- (.10,0);\end{tikzpicture}}

% ms. 12
Pour la propriété universelle de
% ms. 13
\[X\longrightarrow\Pi(X)\ \unsure{\sim|,}\]
% ms. 14
(qui permettra d'écrire
% ms. 15
\[\Pi(X)\simeq X\,\Sigma^{\!-1}\deleted{\unsure{\mathrm{Fl}}},\quad\text{avec }\Sigma=\mathrm{Fl}(X)),\]
% ms. 16–17
autant l'établir tout de suite dans le cadre général où on se donne une partie
% ms. 18
\[\Sigma\subset\mathrm{Fl}(X),\]
% ms. 19
et on définit, pour $\tau\in\mathrm{Ty}$,
% ms. 20
\[\underline{\mathrm{Ch}}_{\Sigma,\tau}(X)\subset\underline{\mathrm{Ch}}_\tau(X)\]
% ms. 21–22
comme la sous-catégorie pleine de $\underline{\mathrm{Ch}}_\tau(X)$ formée des chemins dans $\deleted{\unsure{\Sigma}}X$ dont \uline{\deleted{toutes} les flèches rétrogrades}
\scanpage{327}
\mspage{66}
% archive sheet number: 326
% ms. 1–6
sont dans $X$. On notera que pour qu'un chemin \inserted{$c$} et son opposé $c^{\unsure{\circ|o}}$ soient tous les deux des \uline{$\Sigma$-chemins}, il \unsure{f. et s.} que \uline{toutes} les flèches du chemin $c$ soient dans $\Sigma$ (on \unsure{dira} donc que $c$ est un \uline{$\Sigma$-chemin strict}), puisque les flèches rétrogrades de $c^\circ$ sont les flèches \deleted{\unsure{strictes}} directes de $c$.

% ms. 7–8
Pour que les sous-catégories pleines $\underline{\mathrm{Ch}}_{\tau,\Sigma}(X)$ soient stables pour les foncteurs \unsure{images inverses|image inverse}
% ms. 9
\[f^*:\underline{\mathrm{Ch}}_{\tau'\deleted{\unsure{\Sigma}}}(X)\unsure{\mathord{-}|\longrightarrow}\underline{\mathrm{Ch}}_\tau(X)\]
% ms. 10
associés aux flèches
% ms. 11
\[f:\tau\longrightarrow\tau'\]
% ms. 12–15
dans $\underline{\unsure{\mathrm{Ty}|\mathrm{Tz}}}$, ou simplement \deleted{\unsure{pour}} qu' \deleted{\unsure{elles soient st}}\inserted{il y ait} stabilité pour \deleted{les deux} foncteurs $f^*$ \deleted{qui} \unsure{associés|associé}, \unsure{au} $f:\Delta_-\longrightarrow1$ \deleted{\unsure{(avec $\tau\to1$)}}\inserted{et \deleted{ainsi que pour} $\Delta_-\longrightarrow\Delta_-^2$}, il f. et s. que $\Sigma$ \inserted{soit stable par composition et} contienne toutes les flèches identiques \unsure{:|;}
% ms. 16–17
\[\left\{\begin{aligned}x\in\deleted{\unsure{\mathrm{Ob}}}X&\Longrightarrow\mathrm{id}_x\in\Sigma.\\
u,v\in\Sigma,\ u,v\ \text{\unsure{composables}}&\Longrightarrow uv\in\Sigma\end{aligned}\right.\]
% ms. 18–26
En effet, les flèches dans un chemin $f^*(c')\in\mathrm{Ch}_\tau(X)$ sont soit \deleted{elles} des flèches identiques (qui correspondent aux \deleted{\unsure{couples}} flèches de transition de $I_\tau$ dégénérées pour $f$), soit des flèches parmi celles de $\tau$, ou des composées d'une suite de flèches \deleted{\unsure{compos}}\inserted{consécutives} composables de $\deleted{\unsure{\Sigma}}\tau'$ (qui correspondent aux \deleted{\unsure{couples}}\inserted{flèches} de transition dans $\unsure{I_{\tau'}|\bar I_{\tau'}}$ incluses dans la portée d'une flèche de transition de $I_\tau$), avec conservation des \unsure{signatures}.
\scanpage{328}
\mspage{67}
% archive sheet number: 327
% ms. 1–4
On suppose \unsure{ces|cette} conditions satisfaites, de sorte qu'\unsure{alors} \deleted{\unsure{avec}} les $\mathrm{Ch}_{\tau,\Sigma}^{\unsure{\prime|1}}(X)$ forment les fibres d'une sous-catégorie fibrée \inserted{pleine} de $\underline{\mathrm{Ch}}(X)$ sur $\underline{\unsure{\mathrm{Ty}|\mathrm{Tz}}}$, qu'on dénotera par
% ms. 5
\[\underline{\mathrm{Ch}}_\Sigma(X)\subset\underline{\mathrm{Ch}}(X)\]
% ms. 6–13
Il est évident que cette sous-catégorie est stable par l'opération de composition des chemins\unsure{,|.} (mais non pas \inserted{en général}, on l'a dit, par passage de $c$ à $c^{\unsure{\circ|o}}$ -- (Il faudrait pour cela que $\Sigma$ contienne \uline{toutes} les flèches de $X$, puisque pour \unsure{toute|tout} $u:a\longrightarrow b$, qui est \unsure{tjs} dans $\underline{\mathrm{Ch}}_{\Delta_+,\Sigma}(X)$, $u^\circ$ \deleted{\unsure{$u^\circ$:}}\inserted{\unsure{donné par|décrit par|défini par}} $b\xleftarrow{u}a$ \unsure{dans} $\underline{\mathrm{Ch}}_{\Delta_-\deleted{\unsure{\Sigma}}}(X)$ devrait être un $\Sigma$-chemin, d'où $u\in\Sigma$.

% ms. 14
On va définir une catégorie canonique
% ms. 15
\[\Pi_\Sigma(X)\]
% ms. 16–17
ayant mêmes objets que $X$, et un foncteur canonique
% ms. 18
\[X\longrightarrow\Pi_\Sigma\]
% ms. 19–23
\deleted{transformant} qui est l'identité sur les objets, et qui transforme les $u\in\Sigma$ en flèches inversibles. Pour cela on calque la \deleted{\unsure{défi.}} description du groupoïde fondamental (mais $\Pi_\Sigma(X)$ ne \unsure{sera} pas en général un groupoïde), en posant
% ms. 24
\[\mathrm{Hom}_\Sigma(x,y)=\pi_0(\underline{\mathrm{Ch}}_\Sigma(x,y)),\]
\scanpage{329}
\mspage{68}
% archive sheet number: 328
% ms. 1–2
et en définissant la composition des \deleted{\unsure{classes de}}\inserted{``$\Sigma$-\uline{flèches}''} en termes de la composition des chemins.
% ms. 3–8
L'argument qui prouvait, \deleted{que tout} dans le cas de $\Pi(X)$, que \uline{toute} flèche \inserted{$u$} de $X$, vue comme chemin \inserted{$c$} de type $\Delta_+$, était inversible dans $\Pi(X)$, \deleted{étant} ayant comme inverse $c^{\unsure{\circ|o}}$ de type $\Delta_-$, est encore valable, pourvu que $c^\circ$ soit \inserted{encore} bien un $\Sigma$-chemin, ce qui exige que $u\in\Sigma$. \deleted{\unsure{Il}}

% ms. 9–10
On va l'interpréter en termes d'un foncteur canonique
% ms. 10
\[\deleted{\unsure{p}}k:X\longrightarrow\Pi_\Sigma(X),\tag*{$(\alpha)$}\]
% ms. 10 right
\margin{Sur les flèches, $k$ est défini, \deleted{pour} $\unsure{\forall|\mathrm{tt}}\,u:a\to b$, par $k(u)={}$classe d'homotopie de $\mathrm{ch}(u)$ dans $\underline{\mathrm{Ch}}_\Sigma(a,b)$, où $\mathrm{ch}(u)$ est le chemin de}
% ms. 10 left
% margin note written sideways
\margin{type $\Delta_+$ de $a$ vers $b$ défini par $u$.}
% ms. 11–16
qui se décrit encore comme dans le cas du groupoïde fondamental, en associant à toute flèche \inserted{$u$} de $X$ la classe d'homotopie du chemin $c(u)$ qu'elle définit. On voit donc que ce foncteur canonique transforme les $u\in\Sigma$ en flèches inversibles.

% ms. 17–19
Quand les flèches de $\Sigma$ sont déjà inversibles dans $X$, je dis que le foncteur $(\unsure{\alpha|*})$ est un isomorphisme, \deleted{Pour le voir, \unsure{on va}}\inserted{i.e. que pour tout} $a,b\in\mathrm{Ob}\,X$,
% ms. 20
\[\alpha:\mathrm{Hom}_X(a,b)\longrightarrow\pi_0\,\underline{\mathrm{Ch}}_\Sigma(a,b),\quad u\longmapsto\underset{\substack{\Vert\\\mathrm{ch}(u)}}{\mathrm{ch}_\Sigma(u)}\]
% ms. 21
est bijectif. Pour le voir, on va définir
% ms. 22
\[\beta:\pi_0\bigl(\underline{\mathrm{Ch}}_\Sigma(a,b)\bigr)\longrightarrow\mathrm{Hom}_X(a,b),\]
% ms. 23
\inserted{qui soit} inverse de $\beta$. On définit
\scanpage{330}
\mspage{69}
% archive sheet number: 329
% ms. 1–5
$\beta(c)$, pour une $\Sigma$-flèche de type $\tau$, \inserted{de \unsure{long.|longueur} $n$,} en inversant toutes les flèches rétrogrades (\deleted{\unsure{de}}ce qui est possible, puisqu'elles sont dans $\Sigma$), de façon à obtenir un chemin \inserted{$c^{\unsure{\leftrightarrows|\leftrightsquigarrow}}$} de type $\Delta_+^n$, et en prenant le composé (dans $X$) des flèches de ce chemin $c^{\leftrightarrows}$.
% ms. 5–8
Il est évident donc que $\beta\alpha=\mathrm{id}$, il faut prouver que l'on a aussi $\alpha\beta=\mathrm{id}$, i.e. que dans $\Pi_\Sigma$, les flèches $u=\int c^{\leftrightarrows}$ et $c$ sont égales, i.e. $u$ et $c$ sont homotopes dans $\underline{\mathrm{Ch}}_\Sigma(a,b)$.
% ms. 9
Considérons le chemin $c^{\leftrightarrows}$, de type $\Delta_+^n$ :
% ms. 10
\diagram{\begin{tikzcd}[ampersand replacement=\&,column sep=small]
a=a_0\arrow[r,"u_1"] \& a_1\arrow[r,"u_2"] \& a_2 \& a_{n-1}\arrow[r,"u_n"] \& a_n=b
\end{tikzcd}}
% ms. 11–14
et soit $u_i$ ($i\in[1,n]$) une flèche de ce chemin, correspondant à une flèche de transition \inserted{$\underline{u}_i$} \uline{rétrograde} de $\tau$, donc par construction
% ms. 14
\[u_i=v_{i\deleted{\unsure{u}}}^{-1}\qquad\text{dans } X,\quad\text{avec }v_i\in\Sigma\]
% ms. 15–19
la flèche d'indice \deleted{\unsure{un}} $i$ dans le chemin initial $\tau$ de \deleted{\unsure{$f$}} type $\tau$. \deleted{Or on prouve que \unsure{alors}} \deleted{et cela implique que \unsure{dans les}} Si $\underline{u}_i$ par contre est \deleted{$\mathrm{ch}(u_i)=(\mathrm{ch}(v_i))^{-1}\text{ dans }\Pi_\Sigma(X)$} directe, on a $u_i=v_i$. \deleted{\unsure{Or on a}} (NB On désigne par $v_i$ les $\unsure{i^e|i}$ flèches dans le chemin $c$.) Or on a
% ms. 20
\[\begin{aligned}
&\deleted{\unsure{c(u)=c(u_i)}}\\
&c=c_n*c_{n-1}\cdots c_1\qquad\text{dans }\Pi_\Sigma(X)
\end{aligned}\tag*{$(+)$}\]
% ms. 21–26
et les facteurs $c_i$ ont soit \deleted{donnés par}\inserted{la classe d'homot. de} $\mathrm{ch}(\unsure{v_i|u_i})$ \inserted{$=\mathrm{ch}(u_i)$} (cas où la flèche de transition $\underline{u}_i$ de $\tau$ est directe), soit celle de \deleted{$\mathrm{ch}(u_i)^{\unsure{\circ|o}}$,} du chemin \uline{opposé} au chemin $\mathrm{ch}(v_i):a_i\xrightarrow{v_i}a_{i-1}$ de type $\Delta_+$, donc le chemin $\mathrm{ch}(v_i)^\circ:a_{i-1}\xleftarrow{v_i}a_i$ de type $\Delta_-$. Je dis que
% ms. 27
\[\tikz[baseline=(rel.base)]{\node[inner sep=2pt] (rel) {$\mathrm{ch}(v_i)^\circ\deleted{\unsure{=\mathrm{ch}(v_i)^{-1}}}$};\draw (rel.south west) -- (rel.south east);\draw (rel.south west) -- ++(0,.18);\draw (rel.north west) -- ++(.48,0);}\inserted{=\mathrm{ch}(u_i)}\qquad\text{\uline{dans} }\Pi(X)\tag*{(1)}\]
% ms. 28–30
ce qui prouve que dans $\Pi_\Sigma(X)$, \deleted{le produit} $c$ est égal au composé des $\mathrm{ch}(u_i)$ $(i\in[1,n])$ donc à $c^{\leftrightarrows}$, \deleted{et \unsure{cela} $c^{\leftrightarrows}$ \unsure{est égal à l'objet}} \inserted{et \unsure{il faut|il nous}} prouver de plus que
% ms. 30
\[c^{\leftrightarrows}=\mathrm{ch}(u)\qquad (u=u_nu_{n-1}\cdots u_{\unsure{1|p}})\quad\text{dans }\Pi(X).\tag*{(2)}\]
\scanpage{331}
\mspage{70}
% archive sheet number: 330
% ms. 1–2
Donc il y a deux vérifications à faire (où je suppose à nouveau $\Sigma$ quelconque
% ms. 3
1°) Soit $u\in\Sigma$, \inserted{$u:a\longrightarrow b$,} supposons $\Sigma$ inversible \inserted{dans $X$}. Alors
% ms. 4
\[\mathrm{ch}(u)^{\unsure{\circ|o}}=\mathrm{ch}(u^{-1})\quad\text{dans }\Pi_\Sigma(X)\]
% ms. 5
i.e. les deux chemins de $b$ vers
% ms. 6–7
\diagram{\begin{tikzcd}[ampersand replacement=\&,row sep=small]
b \& a\arrow[l,very thick,"u"']\arrow[l,phantom,"{\deleted{\unsure{\to}}}"{description,very near start}] \& \text{de type }\Delta_-\\
b\arrow[r,"u^{-1}"] \& a \& \text{de type }\Delta_+
\end{tikzcd}}
% ms. 8–10
sont homotopes dans $\underline{\mathrm{Ch}}_\Sigma(b,a)$. Or on sait déjà, \deleted{que} comme $\unsure{k|K}:X\longrightarrow\Pi_\Sigma(X)$ est un foncteur, que
% ms. 11
\[\deleted{\unsure{h|\mathrm{hot}}(\mathrm{ch}(u^{-1}))=(\unsure{h|\mathrm{hot}}(\mathrm{ch}(u))^{-1}\quad\text{dans }\Pi_\Sigma(X)}\]
% ms. 12
\deleted{donc il faut prouver que $\mathrm{ch}(u)^{-1}={}$}
% ms. 13
\[\mathrm{hot}\ \mathrm{ch}(u^{-1})=(\mathrm{hot}\ \mathrm{ch}(u))^{-1}\quad\text{dans }\Pi_\Sigma(X),\]
% ms. 14–18
où $\mathrm{hot}(\bullet)$ désigne la classe d'homotopie d'une \unsure{chaîne} (ici \unsure{chem.|chemin|dans} dans $\underline{\mathrm{Ch}}_\Sigma(b,a)$ \deleted{\unsure{$f$}} et dans $\underline{\mathrm{Ch}}_\Sigma(a,b)$ respectivement. Donc l'énoncé à prouver c'est que $\mathrm{hot}(\mathrm{ch}(u)^\circ)$ est égal à l'\deleted{\unsure{id.|identité|idem}}\inserted{\uline{inverse}} de $\mathrm{hot}(\mathrm{ch}(u))$ dans $\Pi_\Sigma$. C'est donc ce qu'on a prouvé pour \uline{\unsure{tt|$\forall$}} $u\in\mathrm{Fl}(X)$, dans le cas de $\Pi(X)$. L'argument doit marcher, du moment qu'il a un sens \unsure{i.e.|ici,} que $\mathrm{ch}(u)^\circ$ est bel et bien un $\Sigma$-chemin, i.e. $u\in\Sigma$.

% ms. 24–25
2°) La formule (2) signifie que pour tt chemin $c$ de type $\Delta_+^n$ (\unsure{donc} tautologiquement un
\scanpage{332}
\mspage{71}
% archive sheet number: 331
% ms. 1–3
un $\Sigma$-chemin), \inserted{allant de $a$ vers $b$,} sa classe dans $\mathrm{Hom}_\Sigma(a,b)$ n'est autre que celle de $\mathrm{ch}(u)$, où \deleted{\unsure{est}}$\unsure{u}$ est le composé $u_nu_{n-1}\cdots u_1$. Mais on a
% ms. 4
\[c=c_1*c_2*\cdots\unsure{\,.\,|*}c_n\]
% ms. 5
d'où dans $\Pi_\Sigma(X)$
% ms. 6
\[\mathrm{hot}(c)=\mathrm{hot}(c_n)\,\mathrm{hot}(c_{n-1})\cdots\mathrm{hot}(c_1),\]
% ms. 7
or $\mathrm{hot}(c_i)=\mathrm{hot}\,\mathrm{ch}(u_i)=k(u_i)$, donc
% ms. 8–9
\[\begin{aligned}\mathrm{hot}(c)&=k(u_n)k(u_{n-1})\cdots k(u_1)\\
&=k(\underbrace{u_nu_{n-1}\cdots u_1}_{u})\end{aligned}\]
% ms. 10
puisque $k$ est un foncteur, \deleted{\unsure{$=$}} donc
% ms. 11
\[\mathrm{hot}(c)=\mathrm{hot}\,\mathrm{ch}(u),\]
% ms. 12
cqfd.

% ms. 13–16
Notons maintenant que si $(X',\Sigma')$ est un deuxième couple d'une ct. $X'$ et d'un ensemble ``\uline{admissible}'' de flèches de $X'$, et si
% ms. 16
\[f:X\longrightarrow X'\]
% ms. 17
est un foncteur tel que
% ms. 18
\[f(\Sigma)\subset\Sigma',\]
% ms. 19
alors il existe un \uline{unique} foncteur
% ms. 20
\[\bar f=\Pi_\Sigma(f):\deleted{X\Sigma^{-1}}\Pi_\Sigma(X)\longrightarrow\Pi_{\Sigma'}(X')\]
% ms. 21
qui ``prolonge'' $f$, i.e. tel que le carré
\scanpage{333}
\mspage{72}
% archive sheet number: 332
% ms. 1–3
\diagram{\begin{tikzcd}[ampersand replacement=\&]
X\arrow[r,"k_X"]\arrow[d,"f"'] \& \Pi_\Sigma(X)\arrow[d,"{\bar f=\Pi_\Sigma(f)}"]\\
X'\arrow[r,dash,"k_{X'}"] \& \Pi_{\Sigma'}(X')
\end{tikzcd}}
% ms. 4–14
soit commutatif. L'unicité résulte du fait que \deleted{\unsure{l'on} les flèches de $\Pi_\Sigma(X)$ sont}\inserted{sur les objets $\bar f$ doit} être donné par $f$, et que sur les flèches la commutativité \inserted{du carré} implique que $\bar f$ est connu sur les flèches de la forme $k_X(u)$, ainsi \inserted{donc} que sur les inverses de telles flèches, quand ces inverses sont inversibles, donc en particulier sur les $k_X(u)^{-1}$ pour $u\in\Sigma$. Or \deleted{la \unsure{catégo}}\inserted{toute} flèche de $\Pi_\Sigma(X)$, provenant d'un chemin, est \unsure{composé|composée} d'une suite de flèches de cette forme.
% ms. 14–15
Pour l'existence de $\bar f$, on note que $f$ induit des foncteurs
% ms. 16
\[\deleted{\underline{\mathrm{Ch}}_\tau(X)\longrightarrow\underline{\mathrm{Ch}}_\tau}\]
% ms. 17
\[\underline{\mathrm{Ch}}_\tau(X;a,b)\longrightarrow\underline{\mathrm{Ch}}_\tau(X';f(a),f(b))\]
% ms. 18–19
qui commutent aux opérations \inserted{$\unsure{\varphi|g}^*$} sur les chemins associées aux foncteurs
% ms. 20
\[\varphi:\tau\longrightarrow\tau',\]
% ms. 21
d'où un hom de catégories fibrées sur $\underline{\unsure{\mathrm{Ty}|\mathrm{Tz}}}$
% ms. 22
\[f_{a,b}^{\unsure{\varepsilon|\Sigma|\tau}}:\underline{\mathrm{Ch}}(X;a,b)\longrightarrow\underline{\mathrm{Ch}}(X';f(a),f(b)).\]
\scanpage{334}
\mspage{73}
% archive sheet number: 333
% ms. 1
D'où une application sur les $\pi_0$, i.e.
% ms. 2
\[\bar f_{a,b}:\mathrm{Hom}_\Sigma(a,b)\longrightarrow\mathrm{Hom}_\Sigma(f(a),f(b)),\]
% ms. 3
donnée par
% ms. 4
\[\bar f_{a,b}(\mathrm{hot}(c))=\mathrm{hot}(f_{a,b}(c)).\]
% ms. 5–6
On a évidemment, pour deux chemins $\gamma:a\longrightarrow b$ et $\gamma':b\longrightarrow c$
% ms. 6
\[f_{a,\unsure{c|b}}(\gamma'\gamma)=f_{b,c}(\gamma')f_{a,b}(\gamma),\]
% ms. 7–12
ce qui prouve que $\bar f$ est compatible avec la composition. Il est trivial qu'il est compatible avec les identités (car $f_{a,a}$ transforme le chemin vide \inserted{en $a$} en le chemin vide en $f(a)$), donc on a un foncteur $\bar f$, et il rend bel et bien le carré commutatif.

% ms. 13–14
On peut maintenant prouver la propriété universelle de $\Pi_\Sigma(X)$ et de
% ms. 15
\[X\overset{\deleted{\unsure{p}}\unsure{k|p}}{\longleftrightarrow}\Pi_\Sigma(X),\]
% ms. 16–18
savoir que tout foncteur $f:X\longrightarrow X'$ tel que $u\in\Sigma\Longrightarrow f(u)\ \unsure{\mathrm{inv.}|\mathrm{iso.}}$, se factorise \uline{de façon unique} par $\Pi_\Sigma(X)$. L'unicité
% ms. 18–19 left
% slanted
\margin{C'est \unsure{inutile} \unsure{cf|cf.} ci-dessous}
% ms. 19–21
se voit comme tantôt\unsure{$*$|.} Pour l'existence, on note que $f$ \deleted{se factorise par $Y_{\unsure{is|1s}}$, la sous-catégorie pleine de $f$ ayant ses objets, et dont les} \inserted{applique $\Sigma$ dans l'ens. $\Sigma'$\unsure{$*$|:}}
\scanpage{335}
\mspage{74}
% archive sheet number: 334
% ms. 1–2
\deleted{flèches sont les flèches \uline{inversibles} de $Y$. Donc $f$ induit}
% ms. 3
\[\deleted{f_{\unsure{0|o}}:X\longrightarrow{}}\]
% ms. 4–5
des flèches inversibles de $X'$ (lequel ensemble satisfait bien aux deux conditions d'admissibilité).
% ms. 6–7
Donc \deleted{\unsure{par}} ce qui précède, on a un diagramme commutatif
% ms. 8–10
\diagram{\begin{tikzcd}[ampersand replacement=\&,row sep=large,column sep=large]
X\arrow[r,"k"]\arrow[d,"f"'] \& \Pi_\Sigma(X)\arrow[dl,dashed,"\tilde f"]\arrow[d,"{\bar f=\deleted{\unsure{f}}\Pi_\Sigma(f)}"]\\
X'\arrow[r,"{k'}","\sim" description] \& \Pi_{\Sigma'}(X'){,}
\end{tikzcd}}
% ms. 11–15
où $\bar f$ est \unsure{uniquement} déterminé. Comme $k'$ est un isomorphisme (comme on a vu tantôt, puisque $\Sigma'$ \unsure{formé|formée} de fl. \unsure{inv.|inver.}), l'existence et l'unicité de $\deleted{\unsure{\Pi}}\unsure{\tilde f|\bar f}$ \unsure{équivalent|équivaut} à la conclusion annoncée, cqfd.
% ms. between 15 and 16
\figure{\begin{tikzpicture}[x=1cm,y=1cm]\draw (0,.32) -- (.23,0);\end{tikzpicture}}

% ms. 16–19
\uline{Encore un \unsure{sorite}} sur les classes d'homotopie d'éléments d'une catégorie $X$. De façon générale, on voit que \deleted{les \unsure{cl}} $X$ et $\Pi_\Sigma(X)=X\Sigma^{-1}$ ont \unsure{$\Longrightarrow$|mêmes} \deleted{\unsure{classes de}} composantes connexes, en \unsure{ce|si}
\scanpage{336}
\mspage{75}
% archive sheet number: 335
% ms. 1–7
Si $a$ et $b$ \inserted{dans $X$} sont reliés par une flèche de $X\Sigma^{-1}$, i.e. par un $\Sigma$-chemin, alors ils sont dans une même comp. connexe de $X$. Mais l'inverse est \unsure{égalemt} vrai, comme on voit p. ex. en prenant $\Sigma=\mathrm{Fl}(X)$, et en disant que \deleted{\unsure{les composantes connexes du}}\inserted{dans le} \uline{groupoïde} $\Pi(X)$, \deleted{\unsure{si}} deux sommets sont dans une \unsure{m} comp.
% ms. 8–10
ssi $\exists\ x\longrightarrow y$ (relation qui est non seulement transitive et \unsure{reflexive|réflexive}, mais aussi \unsure{symmétrique|symétrique} \deleted{puisque $\Pi(X)$ est}\inserted{dans} un groupoïde). Un chemin
% ms. 11
\[a\xrightarrow{c}b\]
% ms. 12–18
sera appelé aussi une \uline{homotopie} de $a$ vers $b$. On dit que $a$ et $b$ sont \uline{homotopes} s'il existe une \deleted{telle} homotopie \inserted{$c$} de $a$ vers $b$, auquel cas il existe aussi une homotopie \deleted{\unsure{``opp}} ``inverse'' $c^{\unsure{\circ|o}}$ de $b$ vers $a$. Cela signifie aussi que $a$ et $b$ sont dans la $=$ composante connexe de $X$, i.e. \unsure{définissent} le $=$ élément de $\pi_0(X)$.
% ms. 19–26
L'homotopie sera dite \uline{simple} si \unsure{$c$|elle} est de type $\Delta_+$ ou $\Delta_-$, i.e. si elle est donnée par une flèche $u:a\longrightarrow b$ ou $v:b\longrightarrow a$. Si on cherche une homotopie entre $a$ et $b$, on essaiera de la trouver la plus ``courte'' possible. Si on a une homotopie $c$ de longueur $\ell$ de $a$ vers $b$, alors en remplaçant $\tau$ par $\tau_{\mathrm{red}}$, donc $I_\tau$ par $I_{\tau_{\mathrm{red}}}$, la \unsure{sous-échelle|sous-chenille} réduite
\scanpage{337}
\mspage{76}
% archive sheet number: 336
% ms. 1–4
\unsure{correspondante|correspondant}, \deleted{\unsure{on}} on trouve une homotopie de $a$ \unsure{à|et} $b$ de type $\tau_{\mathrm{red}}$, dont la longueur est égale à la ``\uline{longueur utile}'' du type $\tau$ de $c$ (ce qui justifie la terminologie).
% ms. 5
Notons que si
% ms. 6
\[\tau=\Delta_{\varepsilon\deleted{\unsure{+}}}^{\ell_1}\,\Delta_{-\varepsilon\deleted{\unsure{2}}}^{\ell_2}\,\Delta_\varepsilon^{\ell_3}\cdots\Delta_{(-1)^{\unsure{r|n}}\varepsilon}^{\ell_{\deleted{\unsure{n}}r}},\]
% ms. 7
alors
% ms. 8
\[\tau_{\mathrm{red}}=\underbrace{\Delta_\varepsilon\,\Delta_{-\varepsilon}\,\Delta_{+\varepsilon}\cdots\Delta_{(-1)^r\varepsilon}}_{r},\]
% ms. 9–10
où $r$ est la longueur utile de $\tau$ (la longueur \unsure{pas} utile étant $n=\sum_{i\in[1,r]}\ell_i$).

% ms. 11–12
Quand on par\deleted{\unsure{t}}le d'une homotopie entre \uline{foncteurs}
% ms. 12
\[X\longrightarrow Y,\]
% ms. 13–16
il est entendu qu'il s'agit d'homotopies dans la catégorie $\underline{\mathrm{Hom}}(X\unsure{,|;}Y)$\unsure{.|;} (Ou, suivant le contexte, dans \deleted{\unsure{telle}}\unsure{une} sous-catégorie précisée de $\underline{\mathrm{Hom}}(X,Y)$.)
% ms. below 16
\figure{\begin{tikzpicture}[x=1cm,y=1cm]\draw[line width=.55pt] (0,0) .. controls (.18,-.26) and (.62,-.67) .. (.66,-.91);\end{tikzpicture}}
